\section{Estimación por intervalo avanzada y tamaño de muestra}

En las secciones previas introducimos de manera intuitiva el concepto de intervalo de confianza y evaluamos estimaciones por intervalo básicas para la media poblacional. Sin embargo, en el análisis estadístico real y en los proyectos de ciencia de datos, rara vez nos limitamos a estudiar una sola población o a contar con parámetros poblacionales conocidos. Con frecuencia debemos comparar dos tratamientos (como en las pruebas A/B de páginas web), evaluar el comportamiento de proporciones o tasas de conversión, analizar la variabilidad o dispersión de los datos y, de manera crítica ante cualquier estudio, calcular con precisión el tamaño de muestra requerido para alcanzar una potencia o precisión determinada.

En esta sección desarrollamos de forma formal, completa y rigurosa la teoría de estimación por intervalo para diferencias de medias, errores estándares, proporciones, varianzas y diseño del tamaño muestral.

\subsection{Estimación de la media y diferencia entre dos medias}

Cuando se estudian dos poblaciones, el interés principal recae en estimar la diferencia de sus medias, $\mu_1 - \mu_2$, a partir de muestras independientes o pareadas extraídas de cada población. El estimador puntual natural y de mínima varianza para esta diferencia es la diferencia de las medias muestrales: $\bar{X}_1 - \bar{X}_2$.

La construcción del intervalo de confianza para $\mu_1 - \mu_2$ depende críticamente del tipo de diseño de muestreo y del conocimiento o suposición sobre las varianzas poblacionales ($\sigma_1^2, \sigma_2^2$).

\begin{teorema}[Caso 1: Muestras independientes con varianzas conocidas]
	Sean $X_{11}, X_{12}, \ldots, X_{1n_1}$ y $X_{21}, X_{22}, \ldots, X_{2n_2}$ dos muestras aleatorias independientes de tamaños $n_1$ y $n_2$ procedentes de poblaciones con medias $\mu_1, \mu_2$ y varianzas conocidas $\sigma_1^2, \sigma_2^2$. Si las poblaciones son normales o los tamaños muestrales son grandes ($n_1, n_2 \geq 30$ por el Teorema del Límite Central), un intervalo de confianza del $(1-\alpha)100\%$ para la diferencia de medias es:
	\begin{align}
		(\bar{X}_1 - \bar{X}_2) - Z_{\alpha/2} \sqrt{\frac{\sigma_1^2}{n_1} + \frac{\sigma_2^2}{n_2}} \leq \mu_1 - \mu_2 \leq (\bar{X}_1 - \bar{X}_2) + Z_{\alpha/2} \sqrt{\frac{\sigma_1^2}{n_1} + \frac{\sigma_2^2}{n_2}},
	\end{align}
	donde $Z_{\alpha/2}$ es el valor crítico de la distribución normal estándar que deja una probabilidad superior de $\alpha/2$.
\end{teorema}

\begin{teorema}[Caso 2: Muestras independientes con varianzas desconocidas pero iguales]
	Si las varianzas poblacionales son desconocidas pero se puede asumir homoscedasticidad ($\sigma_1^2 = \sigma_2^2 = \sigma^2$), combinamos (\emph{pool}) la información de ambas muestras para obtener una estimación ponderada de la varianza común, denotada por la varianza combinada $S_p^2$:
	\begin{align}
		S_p^2 = \frac{(n_1 - 1)S_1^2 + (n_2 - 1)S_2^2}{n_1 + n_2 - 2},
		\label{eq:varianza_combinada}
	\end{align}
	donde $S_1^2$ y $S_2^2$ son las varianzas muestrales insesgadas respectivas. 
	
	El estadístico estandarizado sigue una distribución $t$ de Student con $\nu = n_1 + n_2 - 2$ grados de libertad. El intervalo de confianza del $(1-\alpha)100\%$ para $\mu_1 - \mu_2$ es:
	\begin{align}
		(\bar{X}_1 - \bar{X}_2) \pm t_{\alpha/2, \, n_1+n_2-2} \, S_p \sqrt{\frac{1}{n_1} + \frac{1}{n_2}}.
	\end{align}
\end{teorema}

\begin{teorema}[Caso 3: Muestras independientes con varianzas desconocidas y diferentes (Welch)]
	Cuando las varianzas poblacionales son desconocidas y no es razonable suponer que sean iguales ($\sigma_1^2 \neq \sigma_2^2$), no podemos utilizar la varianza combinada. En este caso (conocido como el problema de Behrens-Fisher), utilizamos las varianzas muestrales individuales y aproximamos la distribución del estadístico mediante la distribución $t$ de Student con los grados de libertad efectivos de Welch-Satterthwaite:
	\begin{align}
		\nu = \frac{\left( \frac{S_1^2}{n_1} + \frac{S_2^2}{n_2} \right)^2}{\frac{\left(S_1^2 / n_1\right)^2}{n_1 - 1} + \frac{\left(S_2^2 / n_2\right)^2}{n_2 - 1}}.
		\label{eq:df_welch}
	\end{align}
	Si $\nu$ no es un entero, se suele redondear hacia abajo al entero inmediato para mantener una actitud conservadora en la cobertura del intervalo. El intervalo del $(1-\alpha)100\%$ resulta en:
	\begin{align}
		(\bar{X}_1 - \bar{X}_2) \pm t_{\alpha/2, \, \nu} \sqrt{\frac{S_1^2}{n_1} + \frac{S_2^2}{n_2}}.
	\end{align}
\end{teorema}

\begin{definicion}[Caso 4: Muestreo pareado o dependiente]
	Cuando las dos muestras están correlacionadas par a par (por ejemplo, mediciones del mismo paciente o del mismo sistema antes y después de un tratamiento), no podemos tratarlas como independientes. En su lugar, reducimos el problema de dos variables a un problema de una sola variable calculando las diferencias individuales:
	\begin{align}
		D_i = X_{1i} - X_{2i}, \quad \text{para } i = 1, 2, \ldots, n.
	\end{align}
	
	Calculamos la media de las diferencias $\bar{D} = \frac{1}{n}\sum D_i$ y la desviación estándar de las diferencias $S_D = \sqrt{\frac{1}{n-1}\sum (D_i - \bar{D})^2}$. El intervalo de confianza del $(1-\alpha)100\%$ para la diferencia media poblacional $\mu_D = \mu_1 - \mu_2$ es:
	\begin{align}
		\bar{D} - t_{\alpha/2, \, n-1} \frac{S_D}{\sqrt{n}} \leq \mu_D \leq \bar{D} + t_{\alpha/2, \, n-1} \frac{S_D}{\sqrt{n}}.
	\end{align}
\end{definicion}

\subsection{Errores estándares de los principales estimadores}

El \emph{Error Estándar} ($SE$, de \emph{Standard Error}) de cualquier estimador puntual $\hat{\theta}$ es la desviación estándar de la distribución muestral de dicho estimador: $SE(\hat{\theta}) = \sqrt{\text{Var}(\hat{\theta})}$. Constituye el bloque de construcción fundamental del margen de error en todo intervalo de confianza.

Para referencia sistemática en ciencia de datos, la Tabla \ref{tab:errores_estandar} resume los errores estándares teóricos y sus estimadores muestrales empíricos para los parámetros más relevantes en inferencia estadística.

\begin{table}[h!]
	\centering
	\begin{tabular}{lccc}
		\hline
		\textbf{Parámetro ($\theta$)} & \textbf{Estimador ($\hat{\theta}$)} & \textbf{Error Estándar Teórico} & \textbf{Estimación del Error Estándar ($\widehat{SE}$)} \\ \hline
		Media una pob. ($\mu$) & $\bar{X}$ & $\frac{\sigma}{\sqrt{n}}$ & $\frac{s}{\sqrt{n}}$ \\ [1.5ex]
		Dif. medias indep. ($\mu_1 - \mu_2$) & $\bar{X}_1 - \bar{X}_2$ & $\sqrt{\frac{\sigma_1^2}{n_1} + \frac{\sigma_2^2}{n_2}}$ & $\sqrt{\frac{S_1^2}{n_1} + \frac{S_2^2}{n_2}}$ \text{ o } $S_p \sqrt{\frac{1}{n_1} + \frac{1}{n_2}}$ \\ [1.5ex]
		Dif. medias pareadas ($\mu_D$) & $\bar{D}$ & $\frac{\sigma_D}{\sqrt{n}}$ & $\frac{S_D}{\sqrt{n}}$ \\ [1.5ex]
		Proporción pobl. ($p$) & $\hat{p} = \frac{X}{n}$ & $\sqrt{\frac{p(1-p)}{n}}$ & $\sqrt{\frac{\hat{p}(1-\hat{p})}{n}}$ \\ [1.5ex]
		Dif. proporciones ($p_1 - p_2$) & $\hat{p}_1 - \hat{p}_2$ & $\sqrt{\frac{p_1(1-p_1)}{n_1} + \frac{p_2(1-p_2)}{n_2}}$ & $\sqrt{\frac{\hat{p}_1(1-\hat{p}_1)}{n_1} + \frac{\hat{p}_2(1-\hat{p}_2)}{n_2}}$ \\ \hline
	\end{tabular}
	\caption{Resumen de errores estándares para los principales estimadores en estadística inferencial.}
	\label{tab:errores_estandar}
\end{table}

\subsection{Estimación de una proporción y diferencia entre dos proporciones}

El estudio de proporciones poblacionales es vital en el análisis de variables categóricas binomiales o de Bernoulli, siendo la base de las métricas de clasificación, tasas de clics en comercio electrónico y porcentajes de preferencia.

\begin{teorema}[Intervalo de confianza para una proporción (Aproximación normal de Wald)]
	Sea $X$ el número de éxitos en una muestra aleatoria de tamaño $n$ de una población Bernoulli con verdadera proporción $p$. El estimador puntual es la proporción muestral $\hat{p} = X/n$. Por el Teorema del Límite Central, si las condiciones de muestra grande se cumplen ($n\hat{p} \geq 5$ y $n(1-\hat{p}) \geq 5$, o de manera más conservadora $\geq 10$), la distribución muestral de $\hat{p}$ es aproximadamente normal.
	
	El intervalo de confianza del $(1-\alpha)100\%$ está dado por:
	\begin{align}
		\hat{p} - Z_{\alpha/2} \sqrt{\frac{\hat{p}(1-\hat{p})}{n}} \leq p \leq \hat{p} + Z_{\alpha/2} \sqrt{\frac{\hat{p}(1-\hat{p})}{n}}.
	\end{align}
\end{teorema}

\begin{observacion}
	Para muestras pequeñas o moderadas, el intervalo de Wald puede presentar coberturas reales notablemente inferiores al nivel $(1-\alpha)$ nominal debido a la asimetría discreta de la distribución binomial. En tales escenarios, los científicos de datos prefieren el \emph{Intervalo de Wilson (o Score)} o la aproximación de Agresti-Coull (añadir 2 éxitos y 2 fracasos ficticios a la muestra, usando $\tilde{n} = n + 4$ y $\tilde{p} = (X+2)/(n+4)$).
\end{observacion}

\begin{teorema}[Intervalo de confianza para la diferencia de dos proporciones]
	Para comparar las tasas de éxito $p_1$ y $p_2$ de dos poblaciones independientes, extraemos muestras de tamaños $n_1$ y $n_2$ y calculamos las proporciones muestrales $\hat{p}_1 = X_1/n_1$ y $\hat{p}_2 = X_2/n_2$. Si ambas muestras son suficientemente grandes ($n_i\hat{p}_i \geq 5$ y $n_i(1-\hat{p}_i) \geq 5$ para $i=1,2$), el intervalo de confianza de nivel $(1-\alpha)100\%$ para $p_1 - p_2$ es:
	\begin{align}
		(\hat{p}_1 - \hat{p}_2) \pm Z_{\alpha/2} \sqrt{\frac{\hat{p}_1(1-\hat{p}_1)}{n_1} + \frac{\hat{p}_2(1-\hat{p}_2)}{n_2}}.
	\end{align}
\end{teorema}

\subsection{Estimación de la varianza y razón de dos varianzas}

En muchas aplicaciones estadísticas, la variabilidad de un proceso es tan trascendental como su promedio. Por ejemplo, en el control de calidad, en finanzas cuantitativas (donde la varianza mide el riesgo o volatilidad del activo) o al verificar los supuestos del ANOVA y la regresión lineal.

\begin{teorema}[Intervalo de confianza para una varianza poblacional $\sigma^2$]
	Sea $X_1, X_2, \ldots, X_n$ una muestra aleatoria de tamaño $n$ de una población distribuida normalmente con media $\mu$ y varianza $\sigma^2$. El estadístico muestral:
	\begin{align}
		\chi^2 = \frac{(n-1)S^2}{\sigma^2}
	\end{align}
	sigue una distribución chi-cuadrada ($\chi^2$) con $n-1$ grados de libertad.
	
	Por lo tanto, un intervalo de confianza del $(1-\alpha)100\%$ para la varianza poblacional $\sigma^2$ (y sacando raíz cuadrada para la desviación estándar $\sigma$) se obtiene despejando los cuantiles de la distribución $\chi^2$:
	\begin{align}
		\frac{(n-1)S^2}{\chi^2_{\alpha/2, \, n-1}} \leq \sigma^2 \leq \frac{(n-1)S^2}{\chi^2_{1-\alpha/2, \, n-1}},
	\end{align}
	donde $\chi^2_{\alpha/2, n-1}$ y $\chi^2_{1-\alpha/2, n-1}$ son los cuantiles superior e inferior de la distribución chi-cuadrada que dejan un área en las colas derechas de $\alpha/2$ y $1-\alpha/2$, respectivamente.
\end{teorema}

\begin{teorema}[Intervalo de confianza para la razón de dos varianzas $\sigma_1^2 / \sigma_2^2$]
	Sean dos muestras aleatorias independientes de tamaños $n_1$ y $n_2$ provenientes de poblaciones normales independientes con varianzas $\sigma_1^2$ y $\sigma_2^2$. El cociente de las varianzas muestrales estandarizadas:
	\begin{align}
		F = \frac{S_1^2 / \sigma_1^2}{S_2^2 / \sigma_2^2} = \left(\frac{S_1^2}{S_2^2}\right)\left(\frac{\sigma_2^2}{\sigma_1^2}\right)
	\end{align}
	sigue una distribución $F$ de Snedecor con $\nu_1 = n_1 - 1$ grados de libertad en el numerador y $\nu_2 = n_2 - 1$ grados de libertad en el denominador.
	
	El intervalo de confianza del $(1-\alpha)100\%$ para el cociente de varianzas poblacionales $\frac{\sigma_1^2}{\sigma_2^2}$ es:
	\begin{align}
		\frac{S_1^2}{S_2^2} \frac{1}{F_{\alpha/2, \, n_1-1, \, n_2-1}} \leq \frac{\sigma_1^2}{\sigma_2^2} \leq \frac{S_1^2}{S_2^2} \frac{1}{F_{1-\alpha/2, \, n_1-1, \, n_2-1}},
	\end{align}
	donde $F_{\alpha/2, \nu_1, \nu_2}$ es el cuantil superior de la distribución $F$. 
	
	Recordando la propiedad de reciprocidad de la distribución $F$, se cumple que $F_{1-\alpha/2, \nu_1, \nu_2} = \frac{1}{F_{\alpha/2, \nu_2, \nu_1}}$, lo que permite simplificar cálculos cuando las tablas estadísticas solo reportan las colas superiores.
\end{teorema}

\subsection{Estimación del tamaño de una muestra}

Uno de los interrogantes de mayor impacto económico y computacional en la ciencia de datos experimental es: \emph{¿cuántos datos necesitamos?} Tomar una muestra demasiado pequeña puede hacer que el estudio carezca de poder para detectar efectos reales, mientras que una muestra excesivamente grande desperdicia recursos financieros, tiempo e infraestructura computacional.

Para determinar el tamaño muestral $n$ adecuado antes de recolectar los datos, debemos especificar tres elementos:
\begin{enumerate}
	\item El \textbf{nivel de confianza nominal $(1-\alpha)$}, que determina el valor crítico ($Z_{\alpha/2}$ o $t_{\alpha/2}$).
	\item El \textbf{margen de error máximo tolerable ($E$)}, que representa la mitad de la longitud deseada del intervalo de confianza.
	\item Una \textbf{estimación previa de la dispersión o variabilidad poblacional} ($\sigma$ para medias, o una proporción esperada $p^*$ para proporciones).
\end{enumerate}

\begin{teorema}[Tamaño de muestra para estimar una media poblacional $\mu$]
	Al estimar la media de una población normal o bajo el Teorema del Límite Central, el margen de error del intervalo de confianza es $E = Z_{\alpha/2} \frac{\sigma}{\sqrt{n}}$. Despejando $n$, obtenemos el tamaño de muestra requerido:
	\begin{align}
		n = \left( \frac{Z_{\alpha/2} \cdot \sigma}{E} \right)^2.
		\label{eq:n_media}
	\end{align}
	Si la varianza poblacional $\sigma^2$ se desconoce, se puede obtener una estimación preliminar mediante: (a) un estudio piloto previo con una muestra pequeña ($n_0 \approx 30$) para usar la desviación estándar muestral $s$; o (b) la regla empírica del rango ($\sigma \approx \frac{\text{Rango estimado}}{4}$ si la población es aproximadamente normal). Siempre que el resultado no sea un número entero, el valor calculado se redondea \emph{hacia el entero superior siguiente} para garantizar que el margen de error final sea menor o igual a $E$.
\end{teorema}

\begin{teorema}[Tamaño de muestra para estimar una proporción poblacional $p$]
	El margen de error para una proporción poblacional utilizando la aproximación normal es $E = Z_{\alpha/2} \sqrt{\frac{p(1-p)}{n}}$. Despejando $n$, obtenemos:
	\begin{align}
		n = \left( \frac{Z_{\alpha/2}}{E} \right)^2 p^* (1 - p^*),
		\label{eq:n_proporcion}
	\end{align}
	donde $p^*$ es una estimación a priori de la proporción real.
	
	\textbf{Caso conservador (sin información previa):} Si no se dispone de ningún dato preliminar sobre el valor probable de $p$, debemos protegernos ante el peor de los casos en cuanto a variabilidad. La función cuadrática $g(p) = p(1-p)$ alcanza su valor máximo absoluto en $p = 0.5$, donde $g(0.5) = 0.5 \times 0.5 = 0.25$. Sustituyendo este máximo en la ecuación anterior, obtenemos la fórmula más segura y conservadora:
	\begin{align}
		n_{\text{max}} = \frac{Z_{\alpha/2}^2}{4 E^2}.
		\label{eq:n_conservador}
	\end{align}
\end{teorema}

\begin{observacion}
	Cuando la población total de estudio es de tamaño finito y relativamente pequeño (denotado por $N$), y la muestra calculada por las fórmulas anteriores representa más del $5\%$ de la población ($n / N > 0.05$), debemos aplicar el \emph{Factor de Corrección por Población Finita} (FCPF). El tamaño muestral ajustado $n_a$ es:
	\begin{align}
		n_a = \frac{n}{1 + \frac{n-1}{N}}.
	\end{align}
	Este ajuste reduce la cantidad de observaciones requeridas sin sacrificar la precisión estipulada.
\end{observacion}
